\documentclass[11pt]{article}
\usepackage[margin=1in]{geometry}
\usepackage{amsmath,amssymb,amsthm}
\usepackage{booktabs}
\usepackage{graphicx}
\usepackage{microtype}
\usepackage{xcolor}
\usepackage[colorlinks=true,linkcolor=blue!60!black,citecolor=blue!60!black,urlcolor=blue!60!black]{hyperref}
\usepackage[numbers,sort&compress]{natbib}

\newtheorem{theorem}{Theorem}
\newtheorem{proposition}{Proposition}
\newtheorem{corollary}{Corollary}
\newtheorem{remark}{Remark}
\newcommand{\goo}{g_{00}}
\newcommand{\gto}{g_{10}}
\newcommand{\got}{g_{01}}
\newcommand{\gtt}{g_{11}}
\newcommand{\gpi}{g_{\pi}}
\newcommand{\todoref}[1]{\textcolor{red}{[REF: #1]}}

\title{\bf Stale Rollouts Can Pick the Wrong Prompts}
\author{Anonymous\thanks{Draft v0.1 (2026-08-13). Reframed structure following the
pre-registered ``structural absence'' branch; numbers are from completed runs
(v1 gate + v2 corrected, 2 seeds); pending seed harvests will update
Sections~\ref{sec:map} and \ref{sec:diagnosis} only.}}
\date{}

\begin{document}
\maketitle

\begin{abstract}
Reinforcement learning with verifiable rewards (RLVR) increasingly ranks
training prompts with \emph{stale} rollouts---responses generated by an earlier
policy---after a one-sided importance correction: prefix-ratio methods restore
where the current policy goes, future-trace methods restore how it finishes,
and token-ratio influence scores restore neither. We ask when such stale
selection can be trusted, and answer in three tiers.
\textbf{(i)~Possibility (proved).} We organize existing estimators into a
$2{\times}2$ decomposition over prefix occupancy and continuation outcome and
prove that any one-sided correction can reverse the sign of a prompt's policy
gradient at arbitrarily small policy divergence; the failure survives group
normalization, and KL, ESS, and pointwise-cosine dashboards provably cannot
exclude it. Cell \emph{disagreement} is the one stale-computable warning sign:
in a seeded random search over $5{\times}10^4$ two-token environments, every
double reversal was accompanied by disagreement between the two one-sided
estimators.
\textbf{(ii)~Realization map (measured).} Across Qwen2.5 7B/14B and three math
tasks under controlled LoRA drift, whether reversal materializes depends on the
signal regime, which we measure by a corrected split-half oracle floor. The
only regime with measurable ranking signal (floor $0.32$ vs.\ chance $0.10$)
shows one-sided precision loss, while in weak-signal regimes stale scores track
the oracle (Fisher-combined exact $p=1.0{\times}10^{-6}$), and on a saturated
task the selection problem itself degenerates (floor $0.080$, below chance).
\textbf{(iii)~Impossibility (proved + measured).} Two impossibility results
close the loop. First, certifying the top-$k$ decision is out of reach:
alignment margins at the ranking boundary ($0.00$--$0.24^\circ$) sit
\emph{5--8 orders of magnitude in rollout budget} below achievable confidence
radii, across an $8\times$ drift range and all selection fractions 5--25\%.
Second, we prove a \emph{measurement ceiling}: the precision of \emph{any}
estimator measured against a $K$-rollout oracle is bounded by a function of the
split-half floor, so when the floor is near chance no method can measurably
win---our own stale baseline already sits at this ceiling ($0.294$ observed
vs.\ ${\approx}0.30$ ceiling). We therefore contribute a GPU-free diagnosis
procedure that decides, from existing artifacts alone, whether a prompt pool
supports selection at all (conditional-subset precheck plus a $K$-scaling floor
curve with Spearman--Brown extrapolation), and report its pre-registered
verdict on our pools: structural absence. Along the way we document four
evaluation pitfalls that our own pipeline caught---an oracle-sample reuse leak,
shared-jitter tie inflation, cell-budget imbalance, and saturation-shared
overlap inflation---each with a guard. The practical prescription is
\emph{measure the floor first}: selection, and even the evaluation of
selection, is only meaningful where recoverable ranking signal exists, and
today's dashboards do not tell you.
\end{abstract}

\section{Introduction}
\label{sec:intro}

RLVR pipelines score each candidate prompt by ``how much would training on it
help the current policy $\pi$'' and keep the top $k$. Exact scores require
fresh rollouts from $\pi$ for every candidate at every selection round, which
is expensive; practice therefore reuses rollouts and rewards produced by an
earlier behavior policy $\beta$, together with an importance correction. The
correction has two axes: (a)~how the current policy \emph{reaches} a token
(prefix occupancy) and (b)~how the current policy \emph{finishes} after it
(continuation outcome). Existing estimators correct exactly one axis---not by
oversight, but because correcting both multiplies hundreds of token ratios and
explodes the variance. The implicit bet behind this truncation is that the
remaining half induces a small bias that roughly preserves \emph{ranking}.

This paper audits that bet end to end and reports that the situation is worse,
and stranger, than the bias--variance framing suggests. Our claims come in
three tiers, in decreasing order of generality and increasing order of
empirical content:

\begin{enumerate}
\item \textbf{Possibility (proved).} One-sided corrections can reverse the sign
of a prompt's gradient at arbitrarily small policy divergence
(Section~\ref{sec:theory}). Group normalization does not rescue the sign, and
the dashboards practitioners watch---KL bounds, ESS, pointwise gradient
cosine---provably cannot exclude the failure. The one stale-computable warning
sign is \emph{disagreement between the two one-sided cells}.
\item \textbf{Realization map (measured).} Whether reversal materializes
depends on the task--capability regime (Section~\ref{sec:map}). We chart both
directions with a pre-registered protocol: a regime where one-sided selection
measurably loses precision, a regime where stale selection provably tracks the
oracle, and a regime where the selection problem degenerates outright.
\item \textbf{Impossibility (proved + measured).} Where selection matters most,
neither the decision nor its evaluation can be secured at practical budgets:
boundary certification requires 5--8 orders of magnitude more rollouts than
practical (Section~\ref{sec:cert}), and a measurement ceiling caps what any
estimator can even \emph{demonstrate} against a noisy oracle
(Section~\ref{sec:floor}).
\end{enumerate}

The empirical protagonist of the paper is not an estimator but a measurement:
the \emph{split-half oracle floor}, the top-$k$ agreement between two
independent halves of the oracle's own rollouts. The floor measures how much
prompt-ranking signal is recoverable at a given rollout budget; our measurement
ceiling result (Proposition~\ref{prop:ceiling}) makes it the quantity that
decides whether comparative claims about selection methods are meaningful at
all. Applying our own medicine, we ran a pre-registered, GPU-free diagnosis on
our prompt pools (Section~\ref{sec:diagnosis}) and obtained the verdict
\emph{structural absence}: on saturated or tie-degenerate pools, no increase of
the oracle budget up to $K{=}256$ rollouts per prompt is predicted to create a
signal regime. That verdict, rather than a victory lap for a new estimator, is
the paper's empirical headline---and, we argue, the more useful contribution
for practitioners deciding where their fresh-rollout budget should go.

A final theme runs through the experimental sections. Selection evaluation is
easy to get silently wrong in the optimistic direction: over the course of this
project our own pipeline produced---and then caught---four distinct evaluation
artifacts, each of which had initially manufactured a publishable-looking
positive result (Section~\ref{sec:pitfalls}). We document each mechanism and
its guard, because we believe they are endemic to the growing literature on
rollout reuse and data selection for RLVR.

\section{Related work}
\label{sec:related}

\paragraph{Off-policy influence and prompt selection.}
Token-ratio influence scores rank prompts from offline trajectories with a
single-token correction~\citep{cropi}; in our decomposition this is $\goo$
(neither axis restored). Its own appendix already shows the symptom that
motivates our estimand: high pointwise gradient cosine (40/50 prompts above
$0.6$) coexisting with only $28.8\%$ top-10\% neighborhood preservation.
Prefix-ratio corrections~\citep{stepback,cumulative} restore occupancy
($\gto$); their guarantees assume a target-policy advantage that practical
critic-free objectives replace by a behavior-generated outcome
reward---precisely the gap our counterexamples exploit. Future-trace
corrections~\citep{multistep} restore the continuation ($\got$) but not
occupancy when used for per-prompt gradient ranking. Full-trajectory
corrections~\citep{ticgrpo} are exact but variance-explosive on long
responses; they serve as the cost anchor, not a practical selector.

\paragraph{Rollout allocation and learning stability.}
Allocation methods~\citep{gradalign,learnalign} decide where to spend fresh
rollouts to maximize training yield, and stale-data stabilizers keep
\emph{learning} well-behaved under reuse; neither certifies, or even measures,
the correctness of the top-$k$ \emph{selection decision}, which is the object
practitioners actually consume. Classical off-policy correction and subset
selection give the statistical backdrop: conditional importance
sampling~\citep{rowland20}, stationary-distribution correction~\citep{liu19},
doubly-robust policy gradients~\citep{huang20}, and information-complexity
lower bounds for top-$k$ identification~\citep{kaufmann13}. Our certification
impossibility (Section~\ref{sec:cert}) is an empirical instance of the latter's
warning: when value gaps at the boundary vanish, identification cost diverges.

\paragraph{Reliability of noisy rankings.}
Our floor and ceiling analysis imports split-half reliability and the
Spearman--Brown prophecy~\citep{spearman1910,brown1910} into RLVR selection
evaluation. To our knowledge, no prior work in this literature reports the
reliability of its own oracle, and Section~\ref{sec:floor} shows several of our
initially-promising comparisons dissolve once it is reported.

\section{Theory: one-sided corrections and what dashboards cannot see}
\label{sec:theory}

\subsection{A $2\times 2$ decomposition of stale-gradient estimators}

For prompt $x$, response $Y=(A_1,\dots,A_T)$ with terminal verifiable reward
$R(Y)$, the current-policy gradient is
$\gpi(x)=\sum_t \mathbb{E}_{H_t,A_t\sim\pi}[Q_t^{\pi}(H_t,A_t)\,z_t]$ with
$z_t=\nabla_\theta\log\pi(A_t\mid H_t)$. With behavior trajectories, token
ratios $r_j=\pi(A_j|H_j)/\beta(A_j|H_j)$, prefix products $P_t=\prod_{j<t}r_j$
and suffix products $S_t=\prod_{j>t}r_j$, four population quantities arise from
the same stale data:
\begin{align*}
\goo&=\textstyle\sum_t\mathbb{E}_\beta[r_t R z_t]
 &&\text{(occupancy }\beta\text{, continuation }\beta\text{)}\\
\gto&=\textstyle\sum_t\mathbb{E}_\beta[P_t r_t R z_t]
 &&\text{(occupancy }\pi\text{, continuation }\beta\text{)}\\
\got&=\textstyle\sum_t\mathbb{E}_\beta[r_t S_t R z_t]
 &&\text{(occupancy }\beta\text{, continuation }\pi\text{)}\\
\gtt&=\textstyle\sum_t\mathbb{E}_\beta[P_t r_t S_t R z_t]=\gpi.
 &&
\end{align*}
The subscripts state which axis is restored. The exact error decompositions
\[
\gto-\gpi=\sum_t\mathbb{E}_{d_\pi,\pi}\!\big[(Q_t^{\beta}-Q_t^{\pi})z_t\big]
\quad\text{(continuation bias)},\qquad
\got-\gpi=\sum_t\big(\mathbb{E}_{d_\beta,\pi}-\mathbb{E}_{d_\pi,\pi}\big)\!\big[Q_t^{\pi}z_t\big]
\quad\text{(occupancy bias)}
\]
show that the two one-sided estimators remove \emph{different} biases; they are
not substitutes.

\subsection{Sign reversal at arbitrarily small divergence}

\begin{theorem}[One-sided sign reversal]
\label{thm:reversal}
For every $\epsilon>0$ there exist two-token environments and policy pairs
$(\pi,\beta)$ with trajectory KL $O(\epsilon^2)$ such that
$\gpi=-\epsilon/2$ while $\goo=\gto=+\epsilon/2$ (continuation counterexample),
and likewise $\goo=\got=+\epsilon/2$ (occupancy counterexample). Group
normalization with any group size $K\ge 2$ preserves the reversal.
\end{theorem}

The two constructions are given in Appendix~\ref{app:cex}; both flip a
Bernoulli success profile so that the policies agree wherever the retained
correction looks and disagree exactly where the discarded half lives.
Both are machine-checked (raw gradients, exact enumeration of group
normalization for $K\in\{2,4,8\}$ and $\epsilon\in\{0.4,0.1,0.01\}$) by a
self-contained verification script shipped with the code.

\begin{corollary}[Dashboard blindness]
\label{cor:blind}
No bound on trajectory KL, no lower bound on ESS, and no norm-blind,
scale-invariant single-estimator diagnostic (e.g.\ pointwise cosine of the
estimator against itself across checkpoints) can exclude sign reversal of a
prompt's selection score: in the constructions above all such quantities are
driven to their benign limits while the ranking direction is reversed. A
gradient-norm floor or a selection-margin bound is necessary; and since
$\epsilon\to 0$ also drives the absolute gradient to zero, normalized cosines
conceal exactly the missing signal floor.
\end{corollary}

\begin{proposition}[Disagreement accompanies double failure]
\label{prop:disagree}
In the continuation counterexample, $\cos(\gto,\got)=-1$: the two one-sided
cells disagree maximally. In a seeded random search over $5\times 10^{4}$
two-token, two-parameter environments we found \emph{zero} instances of a
double reversal (both one-sided estimators reversed against $\gpi$) accompanied
by cell agreement ($\cos(\gto,\got)>0.9$); the identity $\gtt\equiv\gpi$ was
verified on every sample.
\end{proposition}

Proposition~\ref{prop:disagree} is deliberately modest: absence of
disagreement is \emph{not} proven to imply safety in general architectures with
parameter sharing, disagreement cannot certify \emph{which} cell is right, and
its practical resolution is limited by the same margin collapse quantified in
Section~\ref{sec:cert}. We therefore use cell disagreement only as an
\emph{unreliability signal}---a cheap, stale-computable reason to distrust a
ranking---never as a selector. The mechanism is transparent from the error
decompositions: $\gto$ and $\got$ subtract different bias terms from the same
truth, so a large simultaneous failure of both requires the two bias terms to
agree, which the decomposition makes non-generic.

\section{Measuring selection: floor, ceiling, and when comparison is meaningful}
\label{sec:floor}

\paragraph{Estimand.} All experiments target the decision practitioners
consume: the top-$k$ set by $s_i=\cos(\mu_i,v)$, where $\mu_i$ is the prompt's
mean current-policy gradient and $v$ a validation mean gradient. We report
top-$k$ precision against an on-policy oracle, boundary margins, and rollout
budgets. Pointwise cosine averages are diagnostics, never headlines.

\paragraph{The split-half floor.} The oracle itself is a $K$-rollout Monte
Carlo estimate. We therefore report, for every run, the \emph{corrected
split-half floor}: rollouts are split into two disjoint halves (odd/even
micro-groups), each half induces a top-$k$ set (ties broken by
\emph{independent} per-half jitter; Section~\ref{sec:pitfalls} explains why
this correction matters), and the floor is the mean overlap fraction over 20
jitter pairs. Chance is $k/n$. The floor measures recoverable ranking signal
at budget $K$; it is conservative (an estimator can measurably exceed it), and
Section~\ref{sec:pitfalls} shows it is itself vulnerable to tie artifacts
unless computed exactly this way.

\paragraph{A measurement ceiling.} The floor does more than describe the
oracle: it bounds what any comparison against that oracle can show.

\begin{proposition}[Measurement ceiling]
\label{prop:ceiling}
Model prompt scores as exchangeable Gaussians: true value $t_i$, half-oracle
scores $a_i,b_i$ with $\operatorname{corr}(a,b)=\rho_h$, full-$K$ oracle score
$o_i$ with reliability $\rho_f=2\rho_h/(1+\rho_h)$ (Spearman--Brown). Then for
\emph{any} estimator $s$ computed independently of the oracle's sampling noise,
the expected top-$k$ overlap between $s$ and the $K$-rollout oracle is at most
that of the best possible estimator $s=t$, namely
$\mathrm{Ovl}\big(\sqrt{\rho_f}\big)$, where $\mathrm{Ovl}(c)$ is the expected
top-$k$ overlap of two Gaussian scores with correlation $c$. The observable
floor equals $\mathrm{Ovl}(\rho_h)$, so the ceiling is a computable, monotone
function of the observed floor.
\end{proposition}

Table~\ref{tab:ceiling} gives the mapping at our experimental geometry
($n{=}512$, $k{=}51$), validated by simulation (Appendix~\ref{app:ceiling}).
Two consequences organize everything that follows. First, when the floor is
near chance the ceiling collapses toward chance: \emph{no estimator can
measurably win, and comparative evaluation of selection methods is void}.
Second, an estimator observed \emph{at} the ceiling has plausibly exhausted the
measurable signal---further method development cannot show gains without first
raising the floor (more oracle rollouts, or a pool with more signal).

\begin{table}[t]
\centering
\caption{Floor $\to$ ceiling mapping under Proposition~\ref{prop:ceiling}
($n{=}512$, $k{=}51$, chance $0.100$; simulation, 600--800 replicates/cell).
The v2 weak-signal regime (floors $0.137$--$0.176$) caps \emph{any} method's
measurable precision at ${\approx}0.30$--$0.37$; the observed stale baseline
already sits there.}
\label{tab:ceiling}
\begin{tabular}{lccccccc}
\toprule
observed floor & 0.104 & 0.113 & 0.130 & 0.139 & 0.170 & 0.262 & 0.383\\
\midrule
half-reliability $\rho_h$ & 0.02 & 0.05 & 0.10 & 0.12 & 0.20 & 0.40 & 0.60\\
ceiling on measurable precision & 0.170 & 0.219 & 0.278 & 0.300 & 0.372 & 0.516 & 0.640\\
\bottomrule
\end{tabular}
\end{table}

\paragraph{Retention and its domain.} When comparing regimes we report
retention $(\mathrm{prec}-c)/(\mathrm{floor}-c)$ with $c$ the chance level,
declared \emph{undefined} whenever $\mathrm{floor}-c\le 0.02$; saturated cells
otherwise produce sign-flipping or exploding ratios (the 14B GSM8K cell below
is the cautionary case).

\section{The realization map}
\label{sec:map}

\paragraph{Setup.} Qwen2.5-7B/14B-Instruct; GSM8K, MATH-500, and
DAPO-Math-17k; drift between $\beta$ and $\pi$ is controlled by LoRA RFT for
50--400 steps; gradients are JL-projected; oracle uses $K{=}32$ fresh rollouts
per prompt (8 micro-groups $\times$ 4); pools are $n{=}256$ (v1) or $n{=}512$
(v2) with $k{=}10\%$. The v2 protocol was pre-registered after a full audit
(equal per-cell budgets, seeded tie-breaks, independent jitter, manifests);
judgment criteria were fixed before results were read. Two v2 seeds per task
have completed at the time of writing; remaining seeds update numbers, not
verdicts, and the pre-registered analysis will be re-run on arrival.

\subsection{Where signal exists, one-sided selection can lose it}

The v1 gate (7B GSM8K, floor $0.32$---the strongest signal regime we ever
measured) passed its pre-registered criteria: one-sided estimators lose top-10\%
precision against the uncorrected baseline, and at the largest drift (400 LoRA
steps) the prefix-corrected estimator collapses to precision $0.040$ against a
floor of $0.320$ (overlap $1/25$; exact hypergeometric $P(X\le 1)=0.27$, so we
flag this as a directional observation, not an individually significant
below-chance claim). The misranking propagates: selecting by one-sided scores
changes downstream fine-tuned accuracy (base $0.740$: oracle-selected $0.840$
vs.\ one-sided $0.840$/$0.880$ vs.\ random $0.820$)---though in this
saturation-adjacent regime the spread is small, foreshadowing what follows.
On 7B MATH-500 (floor $0.320$) the \emph{ordering flips}: the uncorrected
$\goo$ is worst ($0.160$) and every stale estimator degrades broadly
($0.28$--$0.64$ of available signal); at 14B MATH-500 (floor $0.240$) the loss
concentrates on the suffix axis (signal retention $\got{=}0.44$ vs.\
$\gto{=}1.00$). Which cell fails is regime-dependent; \emph{that some cell
fails while dashboards stay green} is the invariant, exactly as
Theorem~\ref{thm:reversal} licenses and no more.

\subsection{Where signal is weak, stale selection tracks the oracle}

The corrected v2 runs put every regime we have at floor $0.137$--$0.176$
against chance $0.100$ (Table~\ref{tab:v2}). Here stale scores are \emph{not}
distinguishable from the oracle---in fact they track it: on GSM8K the
uncorrected stale baseline achieves precision $0.294$/$0.235$ (seeds 1/2),
upper-tail exact $p=2.8\times10^{-5}$ and $2.1\times10^{-3}$, Fisher-combined
$p=1.0\times10^{-6}$. Crucially, $0.294$ is \emph{at the measurement ceiling}
(${\approx}0.30$ at floor $0.137$; Table~\ref{tab:ceiling}): the stale
estimator has exhausted what the oracle can resolve, so neither its one-sided
competitors' losses nor any method's gains are measurable in this regime. The
DAPO pool is tie-degenerate (boundary margins exactly $0.0$; precision itself
moves by $\pm 0.08$ under tie-break jitter, so we report bands): all stale
estimators sit at chance, and so does everything else.

\begin{table}[t]
\centering
\caption{v2 corrected results ($n{=}512$, $k{=}51$, chance $0.100$; 2 seeds).
Floors are corrected split-half (independent jitter). DAPO precisions are
jitter-sensitive ($\pm0.08$) due to margin-0 tie degeneracy and shown as
bands.}
\label{tab:v2}
\begin{tabular}{llccccc}
\toprule
run & floor & $\goo$ & $\gto$ & $\got$ & $\gtt$ & regime\\
\midrule
GSM8K s1 & 0.137 & 0.294 & 0.275 & 0.196 & 0.196 & weak signal; stale at ceiling\\
GSM8K s2 & 0.176 & 0.235 & 0.196 & 0.216 & 0.216 & weak signal\\
DAPO s1 & 0.176 & \multicolumn{4}{c}{$0.04$--$0.16$ (tie-degenerate)} & margin $0.0$; chance\\
DAPO s2 & 0.176 & \multicolumn{4}{c}{$0.04$--$0.16$ (tie-degenerate)} & margin $0.0$; chance\\
\bottomrule
\end{tabular}
\end{table}

\subsection{Where capability saturates, selection degenerates}

At 14B on GSM8K (${\sim}95\%$ accuracy) the pool contains 132/256 zero-reward-variance
prompts; every estimator's precision equals the chance level
($0.087$--$0.111$ vs.\ $0.098$), and the oracle disagrees \emph{with itself}
below chance (floor $0.080<0.098$). Selection, ranking, and certification all
degenerate together: no method, and no certificate, can matter here. The
practical corollary is that selection pipelines must scale candidate difficulty
with model capability, or silently pay for machinery that cannot change the
outcome.

\section{Certification is out of reach where it matters}
\label{sec:cert}

Even where a point estimate selects well, one may ask for the decision to be
\emph{certified}: spend fresh rollouts near the ranking boundary until
confidence intervals separate the top-$k$ set. We implemented this
(sequential boundary refinement with shared-validation error propagation;
pseudocode in Appendix~\ref{app:cert}) not as a method contribution but as a
measuring instrument, and it returns a structural negative: across every
completed condition (both scales, three tasks, drift 50--400, selection
fractions 5--25\%) certification fails at budgets $1.02$--$1.04\times$ the
uncorrected cost, because top-$k$ boundary margins are $0.00$--$0.24^\circ$
while achievable confidence radii are $51$--$180^\circ$. The gap is not a
constant factor: with radius $\alpha_v(m)=\arcsin(c/\sqrt m)$, reaching the
observed margins requires $\times 4.1{\times}10^{5}$ to
$\times 1.0{\times}10^{8}$ more observations---5--8 orders of magnitude in
budget, robust to $(\epsilon,\delta)$-PAC slack and deepened validation.
Prompt value is intrinsically dense at the selection boundary: \emph{selection
decisions cannot be certified precisely where selection matters}. Fresh
rollouts spent on certification are wasted; Section~\ref{sec:takeaways} argues
they belong in floor measurement instead.

\section{Four evaluation pitfalls our own pipeline caught}
\label{sec:pitfalls}

Every mechanism below manufactured, at some point in this project, a clean
positive result that survived a first review and died under a targeted audit.
We report them as first-class findings; each comes with a guard now built into
our protocol, and we suspect instances of each exist in the current rollout-reuse
literature.

\paragraph{P1: Oracle-sample reuse (leakage).} Our v1 hybrid intervention
(completing the ``missing half'' with fresh continuations) showed a perfect
21/21-condition recovery, mean $+0.46$/$+0.49$ precision---the would-be causal
centerpiece. A code audit found the treated cell reused rollouts that also
defined the oracle; the shared noise, not causation, produced the recovery. A
pre-registered equal-budget decisive experiment with disjoint samples showed
\emph{no consistent recovery} (e.g.\ GSM8K s1: fully-fresh cell $0.38$ \emph{below}
mixed cell $0.46$). \emph{Guard:} strict sample disjointness between treatment
and oracle (odd/even micro-group separation), enforced structurally.

\paragraph{P2: Shared-jitter tie inflation.} Our first floor implementation
broke score ties with a jitter shared across the two halves. On a
tie-degenerate pool (DAPO, margins $0.0$) this inflated the floor from
$0.176$ to $0.804$---which for a week supported a dramatic (and wrong)
``inverse correlation'' headline. \emph{Guard:} independent per-half jitter
streams, 20 seeded pairs, plus reporting precision \emph{bands} under jitter on
tie-heavy pools.

\paragraph{P3: Cell-budget imbalance.} Comparing estimator cells with unequal
effective rollout counts silently favors the better-funded cell; our v1 hybrid
grid had this flaw on top of P1. \emph{Guard:} equal-$K$ redesign of all cell
comparisons, verified from manifests.

\paragraph{P4: Saturation-shared overlap inflation.} An above-chance floor can
be an artifact of \emph{shared liveness} rather than shared ranking: prompts
that saturate in both halves (all-correct or all-wrong) drop to an exactly-tied
score, and the surviving ``live'' sets overlap between halves, producing
top-$k$ agreement with zero true ranking signal. In a synthetic pool with
\emph{no} ranking signal, saturation sharing alone yields floor $0.126$ at
$K'{=}8$, decaying toward chance ($0.110$) at $K'{=}32$---the signature being
a floor that \emph{decreases} as oracle budget grows. Our measured $K$-scaling
curves (Section~\ref{sec:diagnosis}) show exactly this decreasing shape, which
means even our reported weak-signal floors are best read as upper bounds.
\emph{Guard:} always report the floor as a function of $K'$; treat a
non-increasing floor curve as artifact evidence, and report tie mass alongside.

\section{Diagnosing signal absence without GPUs}
\label{sec:diagnosis}

The pitfalls motivate the final contribution: a pre-registered, GPU-free
procedure that decides whether a pool supports selection \emph{before} any new
experiment is funded. It has two stages, both computed from artifacts every
RLVR selection run already produces.

\paragraph{Stage 1: conditional-subset precheck.} Any proposed pool
intervention definable as a filter on existing prompts (e.g.\ ``keep only live
prompts, $0<\text{pass-rate}<1$'') already exists as a subset of completed
runs. Computing the corrected conditional floor on that subset (with bootstrap
CIs) prechecks the intervention at zero GPU cost: if the conditional floor does
not clear $2\times$ chance in a majority of eligible runs, building the new
pool has no evidential basis. On our pools this returned NO-GO, and the
intervention (a live-only GSM8K pool) was cancelled without spending a
GPU-hour.

\paragraph{Stage 2: $K$-scaling floor curve.} From stored micro-group
gradients the floor is \emph{exactly} recomputed at $K'{=}8,16,24,32$, and
extrapolated to $K'{=}64$--$256$ by Spearman--Brown on the half-reliability
$r_1$ mapped through the Gaussian overlap function, with a calibration check at
the largest observed $K'$. The pre-registered verdict rule: if a majority of
eligible runs are predicted to reach $2\times$ chance by $K'\le 256$,
recommend oracle expansion; otherwise declare structural absence.

\begin{table}[t]
\centering
\caption{$K$-scaling floor curves (observed exact recomputation; prediction by
Spearman--Brown + Gaussian overlap). Chance $0.100$. Note the
\emph{decreasing} observed curves---the P4 artifact signature---and the
near-zero half-reliabilities. Verdict (pre-registered): structural absence.}
\label{tab:kcurve}
\begin{tabular}{llcccccc}
\toprule
run & $r_1$ & $K'{=}8$ & $16$ & $24$ & $32$ & pred.\ $128$ & pred.\ $256$\\
\midrule
GSM8K s1 & $0.045$ & 0.263 & 0.208 & 0.201 & 0.182 & 0.288 & 0.391\\
GSM8K s2 & $-0.005$ & 0.220 & 0.182 & 0.170 & 0.156 & 0.094 & 0.094\\
DAPO s1 & $0.009$ & 0.112 & 0.123 & 0.139 & 0.146 & 0.147 & 0.182\\
DAPO s2 & $-0.038$ & 0.119 & 0.148 & 0.158 & 0.158 & 0.094 & 0.094\\
\bottomrule
\end{tabular}
\end{table}

Table~\ref{tab:kcurve} shows the result: one GSM8K seed extrapolates to
$2\times$ chance at $K'{=}128$ from $r_1{=}0.045$; the other is flat at
chance; both DAPO runs never reach it. Majority not met: \emph{structural
absence}. Two honest caveats are pre-registered with the method. The
calibration check shows the Gaussian model under-predicts observed floors at
$K'{=}32$ (e.g.\ $0.094$ predicted vs.\ $0.156$ observed)---consistent with the
P4 artifact inflating the \emph{observed} values, which makes the verdict
conservative in the safe direction. And the single positive extrapolation
rests on an $r_1$ small enough to be within the P4 artifact's reach, so we
treat it as motivation for a cheap re-check when further seeds land, not as
grounds for a GPU expansion.

\section{Practical guidance, limitations, and outlook}
\label{sec:takeaways}

\paragraph{For practitioners.} (1)~\emph{Measure the floor first.} A
48-prompt fresh probe suffices to estimate the split-half floor; below
$2\times$ chance (a threshold to fix \emph{before} looking), retire selection
and spend nothing---random selection is free and provably indistinguishable.
(2)~\emph{Never budget for boundary certification}; the 5--8 order gap is
structural. (3)~\emph{Scale candidate difficulty with capability}, or
selection degenerates regardless of estimator. (4)~\emph{Report floor and
ceiling with any selection comparison}; a claimed improvement above the
measurement ceiling of its own oracle is an artifact by construction.
(5)~\emph{Distrust one-sided rankings when the two one-sided cells disagree};
agreement guarantees nothing, but disagreement is a free warning.

\paragraph{Limitations.} Empirics are Qwen-2.5 (7B/14B), math tasks, and LoRA
drift; two completed v2 seeds (remaining seeds and a small-model
capability--difficulty sweep are running and will be folded in; the verdict
logic is pre-registered). The ceiling proposition assumes an exchangeable
Gaussian score model (simulation-validated at our geometry; the analytic
program of proving it distribution-free is open). The disagreement result is
an existence-plus-search finding, not a general sufficiency theorem. Theory
counterexamples are two-token constructions---deliberately: they show what
\emph{cannot} be excluded by dashboards, not what typically happens; the
realization map is the empirical complement.

\paragraph{Outlook.} We view the negative results as load-bearing: they
delimit where the field's growing investment in stale-rollout selection can
possibly pay off (unsaturated, tie-free, signal-bearing pools---which our
diagnosis identifies for free) and where it cannot. The measurement-first
protocol---floor, ceiling, preregistration, and the four guards---is applicable
to any selection-method paper, including ours.

\begin{thebibliography}{20}
\bibitem[CROPI]{cropi} Data-Efficient RLVR via Off-Policy Influence Guidance.
ACL 2026. \url{https://arxiv.org/abs/2510.26491}
\bibitem[A Step Back]{stepback} A Step Back: Prefix Importance Ratio
Stabilizes Policy Optimization. \url{https://arxiv.org/abs/2601.22718}
\bibitem[Cumulative]{cumulative} Rethinking Importance Sampling in LLM Policy
Optimization: A Cumulative Token Perspective.
\url{https://arxiv.org/abs/2605.07331}
\bibitem[Multi-Step]{multistep} Multi-Step Likelihood-Ratio Correction for
RLVR. \url{https://arxiv.org/abs/2605.20865}
\bibitem[TIC-GRPO]{ticgrpo} TIC-GRPO: Trajectory Importance Correction for
GRPO. \url{https://arxiv.org/abs/2508.02833}
\bibitem[GradAlign]{gradalign} GradAlign. COLM 2026.
\url{https://arxiv.org/abs/2602.21492}
\bibitem[LearnAlign]{learnalign} LearnAlign.
\url{https://arxiv.org/abs/2506.11480}
\bibitem[Rowland et al.]{rowland20} Rowland et al. Conditional Importance
Sampling for Off-Policy Learning. AISTATS 2020.
\bibitem[Liu et al.]{liu19} Liu et al. Off-Policy Policy Gradient with
Stationary Distribution Correction. UAI 2019.
\bibitem[Huang \& Jiang]{huang20} Huang \& Jiang. From Importance Sampling to
Doubly Robust Policy Gradient. ICML 2020.
\bibitem[Kaufmann et al.]{kaufmann13} Kaufmann \& Kalyanakrishnan. Information
Complexity in Bandit Subset Selection. COLT 2013.
\bibitem[Spearman]{spearman1910} Spearman. Correlation calculated from faulty
data. 1910.
\bibitem[Brown]{brown1910} Brown. Some experimental results in the correlation
of mental abilities. 1910.
\end{thebibliography}

\appendix

\section{Counterexample constructions and machine verification}
\label{app:cex}

\paragraph{Continuation counterexample (reverses $\goo,\gto$).} Two-token MDP;
first action $A_1\sim\mathrm{Bernoulli}(q)$, $q{=}1/2$, identical under both
policies; reward $R=A_2$. Behavior success profile
$Q^\beta(A_1{=}0)=\frac12-\epsilon$, $Q^\beta(A_1{=}1)=\frac12+\epsilon$;
current policy swaps the two values. The gradient w.r.t.\ the logit of $q$ is
$\gpi=-\epsilon/2$ while $\goo=\gto=+\epsilon/2$: the prefix already matches
perfectly, so a perfect prefix ratio changes nothing, and the behavior-policy
\emph{ending} drives the sign. Mean suffix KL is $O(\epsilon^2)$; the 1-D
cosine is $-1$ throughout.

\paragraph{Occupancy counterexample (reverses $\goo,\got$).} Behavior visits
$P(A_1{=}0)=\frac12+\epsilon$, current $\frac12-\epsilon$; the second action is
$\mathrm{Bernoulli}(1/2)$ under both, and
$R=A_2$ if $A_1{=}0$, else $1-A_2$. For the second-action logit,
$\gpi=-\epsilon/2$ while $\goo=\got=+\epsilon/2$: no tokens remain after
$A_2$, so the future ratio is exactly $1$, and the stale \emph{visit}
distribution drives the sign. Trajectory KL is again $O(\epsilon^2)$.

\paragraph{Group normalization.} Both constructions have overall success rate
$1/2$ under both policies; binary group standardization multiplies each
direction by the same nonnegative scalar, hence preserves the reversal (exact
enumeration for $K\in\{2,4,8\}$, $\epsilon\in\{0.4,0.1,0.01\}$ in the
verification script, alongside leave-one-out unbiasedness and the confidence
radius bound used in Section~\ref{sec:cert}).

\section{Measurement ceiling: derivation sketch and validation}
\label{app:ceiling}

Under the model of Proposition~\ref{prop:ceiling}, write the full-$K$ oracle
score as $o=\sqrt{\rho_f}\,t+\sqrt{1-\rho_f}\,\eta$ with $\eta$ independent
noise. For any estimator $s$ independent of $\eta$,
$\operatorname{corr}(s,o)=\operatorname{corr}(s,t)\sqrt{\rho_f}\le\sqrt{\rho_f}$,
with equality iff $s$ is (a monotone transform of) $t$. Expected top-$k$
overlap between two jointly Gaussian scores is increasing in their correlation,
which yields the ceiling $\mathrm{Ovl}(\sqrt{\rho_f})$; the observable floor is
$\mathrm{Ovl}(\rho_h)$ by the same function. Simulation at $n{=}512$,
$k{=}51$ (600--800 replicates per cell; independent-jitter tie handling)
produces Table~\ref{tab:ceiling}; the mapping is monotone and smooth, so
inverting the observed floor through it is well-posed. The Gaussian
assumption matters: heavy-tailed score distributions with a stable head can
exceed the Gaussian ceiling estimate, which is precisely the P4 caveat and why
we report the ceiling alongside, not instead of, the $K$-scaling curve.

\section{Sequential boundary certification (measuring instrument)}
\label{app:cert}

Off-policy scores serve as a prior; every candidate receives a small fresh
micro-group; scores carry confidence balls (shared validation error
propagated); fresh rollouts are then spent only on candidates whose balls
straddle the top-$k$ boundary, until
$\min L_{\text{selected}}>\max U_{\text{rest}}$ or the budget is exhausted, in
which case \emph{certification failure is reported, never hidden}. Variants
(scalar CI, $(\epsilon,\delta)$-PAC with $\epsilon\le 0.05$) behave
identically in all our conditions because the binding constraint is the
boundary margin itself (Section~\ref{sec:cert}).

\section{Preregistration and audit trail}
\label{app:prereg}

All v2 judgment criteria, the precheck GO/NO-GO rule, the $K$-scaling verdict
rule, and the reframing decision tree ("structural absence $\Rightarrow$ this
paper's structure") were registered in the project log before the
corresponding results were computed; the four pitfalls each carry a
before/after commit pair. The verification script, manifests, and all
readouts ship with the code release.

\end{document}
